\section{Numerical Validation and Discussion}
\label{sec:experiments}

This section evaluates LLMBO-MO on the fast-charging problem defined in
Section~\ref{sec:problem}. All protocols are evaluated using the
electrochemical--thermal--aging model in Section~\ref{sec:model_section}. We
first compare LLMBO-MO with representative multiobjective optimizers under the
same simulation budget. The available ablation studies are then used to
examine LLM-informed initialization and the regional-guidance branch, followed
by a descriptive comparison across LLM backends and an analysis of
representative Pareto-optimal charging protocols.

\subsection{Experimental Setup}

The Chen2020 parameterization of the LG INR21700-M50 cell~\cite{ref8} is used
as the primary case, while Ecker2015~\cite{ref9} is additionally considered to assess robustness across battery parameterizations.
For each charging protocol, the three objectives are
charging time $t_c$, peak temperature rise $T_p$, and aging indicator $Q_s$.
The main comparison uses six initial evaluations followed by 50 optimization
iterations, giving 56 battery-model evaluations per run. Five matched trials
are used for all methods in the main comparison.

Hypervolume (HV) is adopted as the main multiobjective performance metric.
Because the true Pareto front is unknown, HV is used to assess both convergence
and diversity of the nondominated set; a larger HV indicates a better solution
set under the same reference point. Within each battery parameterization, all
methods use the same objective definitions, simulation budget, and HV
calculation. The battery and algorithm settings are summarized in
Tables~\ref{tab:battery_settings} and~\ref{tab:algorithm_settings}.

\begin{table}[!t]
\caption{Battery and Simulator Settings}
\label{tab:battery_settings}
\centering
\footnotesize
\setlength{\tabcolsep}{3.0pt}
\begin{tabular}{@{}>{\raggedright\arraybackslash}p{0.58\columnwidth}
                    >{\raggedright\arraybackslash}p{0.34\columnwidth}@{}}
\toprule
Parameter & Value \\
\midrule
Primary parameterization & Chen2020, LG INR21700-M50 \\
Additional parameterization & Ecker2015 \\
Chen2020 nominal capacity & 5.0 Ah \\
Chen2020 cell voltage range & 2.50--4.20 V \\
Chen2020 cell format & 21700 cylindrical, 70 $\times$ 21 mm \\
Initial/ambient temperature & 298.15 K \\
Maximum cell temperature & 328.15 K \\
\bottomrule
\end{tabular}
\end{table}

\begin{table}[!t]
\caption{Main Parameter Settings of LLMBO-MO}
\label{tab:algorithm_settings}
\centering
\footnotesize
\setlength{\tabcolsep}{3.0pt}
\begin{tabular}{@{}>{\raggedright\arraybackslash}p{0.62\columnwidth}
                    >{\raggedright\arraybackslash}p{0.30\columnwidth}@{}}
\toprule
Parameter & Value \\
\midrule
Total simulator budget & 56 \\
Number of trials & 5 \\
Initial LLM candidate pool & 15 \\
Selected fraction of LLM candidates & 0.2 \\
Soft-penalty weight $\omega_p$ & \OmegaPenalty \\
Monotonicity weight $\omega_m$ & \OmegaMonotone \\
Diversity weight $\omega_d$ & \OmegaDiversity \\
Preferred terminal SOC span $\Delta s_K^{\mathrm{soft}}$ & \SoftDeltaSK \\
Region-Lift horizon $T_{\mathrm L}$ & \RegionHorizon \\
Number of regional anchors $M$ & \RegionAnchors \\
Confidence gate $c_{\min}$ & \RegionConfidenceThreshold \\
Confidence scale $\alpha_c$ & \RegionConfidenceScale \\
Mean/max shift budgets & \RegionMeanBudget~/~\RegionMaxShift \\
Tchebycheff coefficient $\rho_{\mathrm{aug}}$ & \RhoAug \\
Objective-scale ratio $\rho_R$ & \RhoR \\
\bottomrule
\end{tabular}
\end{table}

Four competitors are considered. ParEGO is the primary BO baseline because it
shares the decomposition--GP structure of the numerical backbone used in
LLMBO-MO~\cite{ref4}. DISK and PIMD are included as recent surrogate-assisted
multiobjective optimizers~\cite{refDISK,refPIMD}, while NSGA-II provides a
classical evolutionary baseline~\cite{ref5}.

\subsection{Performance Comparison}

Figure~\ref{fig:chen_hv_5way} shows the HV trajectories on Chen2020, and
Table~\ref{tab:chen_results} reports the final values. LLMBO-MO reaches higher
HV during the early part of the search and retains the largest mean HV at the
end of the 56-evaluation budget.

\begin{figure}[!t]
\centering
\includegraphics[width=0.99\columnwidth]{chen2020_hv_5way_600dpi.png}
\caption{HV trajectories of the five methods on Chen2020.}
\label{fig:chen_hv_5way}
\end{figure}

\begin{table}[!t]
\caption{Final HV on Chen2020}
\label{tab:chen_results}
\centering
\footnotesize
\setlength{\tabcolsep}{4.0pt}
\begin{tabular}{@{}lc@{}}
\toprule
Method & Final HV (mean $\pm$ SD) $\uparrow$ \\
\midrule
LLMBO-MO & \textbf{\ChenLLMBOMean\,$\pm$\,\ChenLLMBOStd} \\
ParEGO & \ChenParEGOMean\,$\pm$\,\ChenParEGOStd \\
NSGA-II & \ChenNSGAMean\,$\pm$\,\ChenNSGAStd \\
DISK & \ChenDISKMean\,$\pm$\,\ChenDISKStd \\
PIMD & \ChenPIMDMean\,$\pm$\,\ChenPIMDStd \\
\bottomrule
\end{tabular}
\end{table}

After 56 evaluations, LLMBO-MO achieves a mean HV of
\ChenLLMBOMean\,$\pm$\,\ChenLLMBOStd, compared with
\ChenParEGOMean\,$\pm$\,\ChenParEGOStd\ for ParEGO, corresponding to a
\ChenFinalGain\ improvement. Its final mean HV is also higher than those of
NSGA-II, DISK, and PIMD. The advantage appears before the simulation budget is
exhausted, consistent with the intended role of LLM-informed initialization:
battery-domain knowledge is introduced when only a small number of
task-specific observations are available.

To determine whether the advantage over ParEGO is specific to Chen2020, the
comparison is repeated using Ecker2015. Figure~\ref{fig:ecker_hv} and
Table~\ref{tab:ecker_results} show the results.At evaluation 56, they reach
\EckerLLMBOMean\,$\pm$\,\EckerLLMBOStd\ and
\EckerParEGOMean\,$\pm$\,\EckerParEGOStd. The advantage over ParEGO is
therefore also observed under the Ecker2015 parameterization.

\begin{figure}[!t]
\centering
\includegraphics[width=0.99\columnwidth]{ecker2015_hv_600dpi.png}
\caption{HV trajectories of LLMBO-MO and ParEGO on Ecker2015.}
\label{fig:ecker_hv}
\end{figure}

\begin{table}[!t]
\caption{HV Comparison on Ecker2015}
\label{tab:ecker_results}
\centering
\footnotesize
\setlength{\tabcolsep}{4.0pt}
\begin{tabular}{@{}lcc@{}}
\toprule
Method & Evaluation 56 \\
\midrule
ParEGO &
\EckerParEGOMean\,$\pm$\,\EckerParEGOStd \\
LLMBO-MO &
\textbf{\EckerLLMBOMean\,$\pm$\,\EckerLLMBOStd} \\
\bottomrule
\end{tabular}
\end{table}

\subsection{Ablation Study}

The preceding comparison evaluates the complete LLMBO-MO framework. We next
use the ablation studies to determine which part of the LLM assistance is most
clearly supported by the current evidence.

\subsubsection{LLM-Informed Initialization}

The first study compares three initialization strategies: random
initialization, a detailed generic prompt, and a battery-specific prompt. Ten
matched trials are used, and all three variants follow the same numerical BO
procedure after initialization. The final HV values under a 26-evaluation
budget are reported in Table~\ref{tab:prompt_study}.

\begin{table}[!t]
\caption{Effect of Different Initialization Strategies}
\label{tab:prompt_study}
\centering
\footnotesize
\setlength{\tabcolsep}{4.0pt}
\begin{tabular}{@{}lc@{}}
\toprule
Initialization & Final HV (mean $\pm$ SD) $\uparrow$ \\
\midrule
Random & 0.31062 $\pm$ 0.04303 \\
Generic prompt & 0.31371 $\pm$ 0.04675 \\
Battery-specific prompt & \textbf{0.36890 $\pm$ 0.00824} \\
\bottomrule
\end{tabular}
\end{table}

The generic prompt performs similarly to random initialization. In contrast,
the battery-specific prompt increases the mean final HV by \PromptHVGain\ and
substantially reduces the variation among trials. The paired difference
remains significant after Holm correction ($p=\PromptHolmP$). These results
indicate that the gain is associated with battery-specific information in the
prompt rather than with prompting alone.

\subsubsection{Regional Guidance}

The second study compares four variants under the same short budget: the
numerical BO backbone without LLM assistance, LLM-informed initialization
only, regional guidance only, and the complete LLMBO-MO framework. Each run
contains six initial evaluations followed by ten BO iterations.

This experiment evaluates the regional-guidance branch at the system level; it
is not a matched Region-Lift on/off test with otherwise identical regional
preferences. The results are therefore used to assess the practical effect of
the available regional-guidance branch, not to claim an independent causal
gain from the posterior-covariance lift itself.

\begin{figure}[!t]
\centering
\includegraphics[width=0.99\columnwidth]{ablation_hv_box.png}
\caption{Final HV of the four ablation variants on Chen2020.}
\label{fig:ablation_hv_box}
\end{figure}

\begin{table}[!t]
\caption{Ablation Results on Chen2020}
\label{tab:ablation}
\centering
\footnotesize
\setlength{\tabcolsep}{4.0pt}
\begin{tabular}{@{}lc@{}}
\toprule
Variant & Final HV (mean $\pm$ SD) $\uparrow$ \\
\midrule
No LLM & \AblationBaselineMean\,$\pm$\,\AblationBaselineStd \\
+ Initialization & \AblationWarmMean\,$\pm$\,\AblationWarmStd \\
+ Regional guidance & \AblationRegionMean\,$\pm$\,\AblationRegionStd \\
LLMBO-MO & \textbf{\AblationFullMean\,$\pm$\,\AblationFullStd} \\
\bottomrule
\end{tabular}
\end{table}

Initialization provides the main improvement in this short-budget study: the
mean HV increases from \AblationBaselineMean\ to \AblationWarmMean. Regional
guidance alone gives a smaller change, reaching \AblationRegionMean. With both
components enabled, LLMBO-MO obtains the largest mean HV,
\AblationFullMean, and the smallest dispersion among the four variants. The
available evidence therefore identifies LLM-informed initialization as the
more stable source of improvement. A dedicated matched experiment is still
required to isolate the independent contribution of the finalized Region-Lift.

\subsection{Sensitivity to the LLM Backend}
\label{sec:llm_backends}

A supplementary experiment compares three LLM backends under the same
Chen2020 prompt and numerical BO settings. This comparison is intended as a 
descriptive assessment of backend sensitivity rather than a statistical comparison.
\begin{table}[!t]
\caption{HV Obtained with Different LLM Backends}
\label{tab:llm_backends}
\centering
\footnotesize
\setlength{\tabcolsep}{4.5pt}
\begin{tabular}{@{}lccc@{}}
\toprule
Backend & Initial HV & HV at 15 & Final HV \\
\midrule
DeepSeek-V4-Flash & 0.341225 & \textbf{0.373845} & 0.385986 \\
DeepSeek-V4-Pro & 0.345274 & 0.372526 & \textbf{0.387931} \\
GLM-5.3-Flash & \textbf{0.346248} & 0.372527 & 0.385324 \\
\bottomrule
\end{tabular}
\end{table}

The three backends produce similar final HV values, ranging from 0.385324 to
0.387931. This single-seed result does not support a statistical comparison
among backends, but it indicates that the observed behavior is not unique to
one LLM implementation.

\subsection{Analysis of Pareto-Optimal Charging Protocols}

HV summarizes the overall quality of the nondominated set, whereas the
practical output of the optimizer is a collection of charging protocols. We
therefore select three representative solutions from the Chen2020 Pareto set
to examine their objective-space trade-offs and charging trajectories.
Figure~\ref{fig:pareto_profiles} shows their locations in objective space, and
Table~\ref{tab:representative_protocols} lists the corresponding protocol
parameters.

\begin{figure}[!t]
\centering
\includegraphics[width=\columnwidth]{pareto_protocols_ae.pdf}
\caption{Representative Pareto-optimal charging protocols on Chen2020.}
\label{fig:pareto_profiles}
\end{figure}

\begin{table*}[!t]
\caption{Representative Pareto-Optimal Charging Protocols}
\label{tab:representative_protocols}
\centering
\footnotesize
\setlength{\tabcolsep}{5.0pt}
\begin{tabular}{@{}c l r r r@{}}
\toprule
Protocol & $I_1/I_2/I_3/\Delta s_1/\Delta s_2/\Delta s_3$ &
$t_c$ (s) & $T_p$ (K) & $Q_s$ \\
\midrule
A & 6.000/5.000/3.000/0.400/0.300/0.100 & 2880 & 7.567 & 1.261 \\
B & 2.782/2.000/3.000/0.211/0.100/0.489 & 5199 & 3.069 & 0.680 \\
C & 2.000/2.000/2.000/0.245/0.153/0.403 & 7200 & 1.528 & 0.640 \\
\bottomrule
\end{tabular}
\end{table*}

Protocol A represents the fast-charging end of the Pareto set. Its first two
stages use high currents, giving the shortest charging time of 2880~s, but also
the largest temperature rise and aging indicator among the three
representative solutions. Protocol C uses a constant 2~A current and requires
7200~s to reach the target SOC, while giving the smallest temperature rise.
Protocol B lies between these cases and provides an intermediate compromise.
The three solutions illustrate the expected trade-off: increasing charging
current reduces charging time but generally increases thermal and aging costs.

The corresponding charging trajectories are shown in
Fig.~\ref{fig:abc_charging_profiles}. Protocol A reaches the target SOC first
and exhibits the largest temperature increase. Protocol C charges more slowly
and maintains the lowest temperature, while Protocol B gives an intermediate
response. All three protocols remain below the prescribed voltage limit.

\begin{figure}[!t]
\centering
\includegraphics[width=\columnwidth]{abc_charging_profiles.pdf}
\caption{Charging trajectories of protocols A--C on Chen2020.}
\label{fig:abc_charging_profiles}
\end{figure}

Overall, the numerical results support two conclusions. First, LLMBO-MO
achieves the highest final mean HV among the five methods compared on
Chen2020, and its advantage over ParEGO is also observed under Ecker2015.
Second, the ablation studies identify battery-specific initialization as the
clearest empirical source of improvement in the current evidence. The
independent contribution of the finalized Region-Lift should be assessed
through a strictly matched mechanism-level experiment.
